\chapter{Euler's Totient Function and Euler's Theorem}

\section{What the Totient Function Counts}
Euler's totient function, written \(\varphi(n)\), counts how many integers from
\(1\) to \(n\) are coprime to \(n\) (that is, have greatest common divisor \(1\)
with \(n\)).\cite{wiki:euler-totient}

In symbols:
\[
  \varphi(n)=\#\{\,k \in \{1,\dots,n\} : \gcd(k,n)=1\,\}.
\]

Examples:
\begin{itemize}
  \item \(\varphi(1)=1\), because \(\gcd(1,1)=1\).
  \item \(\varphi(9)=6\), because the numbers
        \(1,2,4,5,7,8\) are coprime to \(9\), while \(3,6,9\) are not.
\end{itemize}


\section{Euler's Product Formula}
We can count the number of integers in \(\{1,\dots,n\}\) that are coprime to \(n\)
using the inclusion-exclusion principle: \index{Inclusion--Exclusion Principle}

If \(n\) has distinct prime divisors \(p_1, p_2, \dots, p_r\),
the total number \(\varphi(n)\) is n minus the multiples of \(p_i\), i=1,2,\dots,r,
plus the multiples of \(p_i\) \(p_j\), i,j=1,2,\dots,r, minus the multiples of 
\(p_i\) \(p_j\) \(p_k\), i,j,k=1,2,\dots,r, and so on.
\[
  \varphi(n)
  = n
  - \sum_{i=1}^{r}\frac{n}{p_i}
  + \sum_{1\le i<j\le r}\frac{n}{p_i p_j}
  - \sum_{1\le i<j<k\le r}\frac{n}{p_i p_j p_k}
  + \cdots
  + (-1)^r\frac{n}{p_1p_2\cdots p_r}.
\]

The result can be written as a product of the form:
\[
  \varphi(n)=n\prod_{p\mid n}\left(1-\frac{1}{p}\right).
\]

\subsection*{Example: \texorpdfstring{\(\varphi(20)\)}{phi(20)}}
Since \(20=2^2\cdot 5\),
\[
  \varphi(20)=20\left(1-\frac12\right)\left(1-\frac15\right)
  =20\cdot \frac12\cdot \frac45=8.
\]
So there are 8 integers in \(\{1,\dots,20\}\) that are coprime to \(20\).

\section{Euler's Theorem}
If \(\gcd(a,n)=1\), then
\[
  a^{\varphi(n)} \equiv 1 \pmod n.
\]

\subsection*{Elementary Proof (No Abstract Algebra)}
This proof follows the classic reduced-residue permutation argument (see
\cite{mse:euler-theorem-no-abstract-algebra}).

Let
\[
  R=\{r_1,r_2,\dots,r_{\varphi(n)}\}
\]
be all integers in \(\{1,\dots,n\}\) that are coprime to \(n\).

Multiply each element by \(a\), where \(\gcd(a,n)=1\):
\[
  ar_1,ar_2,\dots,ar_{\varphi(n)}.
\]
Reduce these modulo \(n\).  Then:
\begin{itemize}
  \item each reduced value is still coprime to \(n\), so it is again in \(R\);
  \item no two reduced values are equal modulo \(n\): if
        \(ar_i\equiv ar_j\pmod n\), then \(r_i\equiv r_j\pmod n\), hence
        \(r_i=r_j\).
\end{itemize}

So multiplication by \(a\) permutes the set \(R\).  Therefore their products are
congruent:
\[
  (ar_1)(ar_2)\cdots(ar_{\varphi(n)})
  \equiv
  r_1r_2\cdots r_{\varphi(n)}
  \pmod n.
\]
Hence
\[
  a^{\varphi(n)}(r_1r_2\cdots r_{\varphi(n)})
  \equiv
  r_1r_2\cdots r_{\varphi(n)}
  \pmod n.
\]
Because each \(r_i\) is coprime to \(n\), the product
\(r_1r_2\cdots r_{\varphi(n)}\) is also coprime to \(n\), so we can cancel it.
We get
\[
  a^{\varphi(n)} \equiv 1 \pmod n.
\]

\subsection*{Examples}
\begin{itemize}
  \item \textbf{\(a=2,\ n=5\):} \(\gcd(2,5)=1\) and \(\varphi(5)=4\), so Euler's theorem predicts
        \(2^4\equiv 1 \pmod 5\).  Indeed, \(2^4=16\), and \(16\equiv 1 \pmod 5\).

  \item \textbf{\(a=3,\ n=10\):} \(\gcd(3,10)=1\) and \(\varphi(10)=4\), so
        \(3^4\equiv 1 \pmod{10}\).  Indeed, \(3^4=81\), and
        \(81\equiv 1 \pmod{10}\).

  \item \textbf{\(a=2,\ n=9\):} \(\gcd(2,9)=1\) and \(\varphi(9)=6\), so
        \(2^6\equiv 1 \pmod 9\).  Indeed, \(2^6=64\), and \(64\equiv 1 \pmod 9\).

  \item \textbf{\(a=5,\ n=12\):} \(\gcd(5,12)=1\) and \(\varphi(12)=4\), so
        \(5^4\equiv 1 \pmod{12}\).  Indeed, \(5^2=25\equiv 1 \pmod{12}\), hence
        \(5^4=(5^2)^2\equiv 1^2\equiv 1 \pmod{12}\).
\end{itemize}

The theorem can be used to reduce large powers modulo \(n\).  For example, to
find the ones digit of \(7^{222}\), compute \(7^{222}\pmod{10}\).  Since
\(\gcd(7,10)=1\) and \(\varphi(10)=4\), Euler's theorem gives
\[
  7^4 \equiv 1 \pmod{10}.
\]
Now write \(222 = 4\cdot 55 + 2\).  Then
\[
  7^{222}
  \equiv 7^{4\cdot 55 + 2}
  \equiv (7^4)^{55}\cdot 7^2
  \equiv 1^{55}\cdot 7^2
  \equiv 49
  \equiv 9 \pmod{10}.
\]
So the ones digit of \(7^{222}\) is \(9\).



Euler's theorem is one of the reasons \(\varphi(n)\) is so important in
cryptography: RSA relies on this modular-exponentiation structure.
